% #----------------------------------------------#
% |   Autor Plantilla: Eduardo Cañete Carmona    |
% |   Fecha creación: 01/07/2019                 |
% |   Ultima actualización: 14/03/2026           |
% |   Versión: 3.7                               |
% #----------------------------------------------#

\documentclass[a4paper, 12pt]{book}
\usepackage[spanish,es-tabla]{babel}
\usepackage[utf8]{inputenc}

\usepackage[left = 35mm, right = 25mm, top = 25mm, height = 206mm]{geometry}
\usepackage{fancyhdr}
\usepackage{graphicx}
\usepackage{lipsum}
\usepackage[square,numbers]{natbib}
\usepackage[skip=12pt,font=small]{caption}  %Tamaño de la leyenda
\usepackage{subcaption}
\usepackage{amsmath}
\usepackage{afterpage}
\usepackage{titlesec}
\usepackage{color, colortbl}
\usepackage{listings}
%\usepackage{xcolor}
\usepackage[dvipsnames,svgnames,table]{xcolor}
\usepackage{float}
\usepackage{eurosym}
\usepackage{enumitem}
\usepackage{amssymb}
\usepackage{array}
\usepackage{wrapfig}
\usepackage{multirow}
\usepackage{tabularx}
\usepackage{url}
\usepackage{eurosym}
\usepackage{appendix}
\usepackage{pdfpages}
\usepackage{pdflscape} %poner pdf horizontal
\usepackage[hidelinks]{hyperref}%para el hiperlink de las secciones
\usepackage[newfloat]{minted}
\usepackage[nottoc]{tocbibind}
\usepackage{glossaries}

\usepackage{ulem} %subrayado y tachado
\usepackage{pifont} % símbolos como tick y X
\usepackage{comment} % comentar bloques grandes
%\usepackage{longtable}
\usepackage{seqsplit} % dividir texto largo automáticamente


\usepackage{tikz} % Diagramas de flujo
\usetikzlibrary{shapes.geometric, arrows} %Libería de flechas

% FONTs
\usepackage{fontenc}
\usepackage{textcomp}           % Needed for new symbols like € ?

\usepackage[scaled]{berasans} % Font for the cover similar to Vera 33
%\renewcommand*\familydefault{\sfdefault}  %% To use as the base font of the document is to be sans serif

\usepackage{datetime}
%\newdateformat{monthyeardate}{\monthname, \THEYEAR}
\newcommand\Monthname[1][EMPTY]{
  \ifnum1=#1Enero\else
  \ifnum2=#1Febrero\else
  \ifnum3=#1Marzo\else
  \ifnum4=#1Abril\else
  \ifnum5=#1Mayo\else
  \ifnum6=#1Junio\else
  \ifnum7=#1Julio\else
  \ifnum8=#1Agosto\else
  \ifnum9=#1Septiembre\else
  \ifnum10=#1Octubre\else
  \ifnum11=#1Noviembre\else
  \ifnum12=#1Diciembre\else
  \fi\fi\fi\fi\fi\fi\fi\fi\fi\fi\fi\fi
}
\newdateformat{monthyeardate}{%
  \Monthname[\THEMONTH], \THEYEAR}
%\newdateformat{monthyeardate}{\Monthname[\THEMONTH], \THEYEAR}

% Entorno para código
\newenvironment{code}{\captionsetup{type=listing}}{}    

% Cambiamos el caption de los listing a español
\renewcommand{\listingname}{Código}    

 % Cambiamos el idioma del índice de códigos
\renewcommand*{\listlistingname}{Índice de Códigos}    

\newfloat{diagrama}{htbp}{loc}
\floatname{diagrama}{Diagrama}
\newcommand{\listofdiagrams}{\listof{diagrama}{Índice de Diagramas de Flujo}}
\floatstyle{plaintop}
\restylefloat{table}

\graphicspath{ {Imagenes/} }

\input{_datos_proyecto}
\input{Configuracion/confg_colores}
\input{Configuracion/confg_codigo}
\input{Configuracion/confg_diagrama_flujo}
%----------------------------------------------------------------------------------------
%	ENCABEZADOS Y PIE DE PAGINA
%----------------------------------------------------------------------------------------

\pagestyle{fancy}
	\fancypagestyle{cuerpo}{
      \fancyfoot{}
      \fancyfoot[L]{Trabajo Fin de Grado} %Autor
      \fancyfoot[C]{}
      \fancyfoot[R]{- \thepage \ -}
      \fancyhead{}
      \fancyhead[L]{\leftmark} %Nombre del documento
      \fancyhead[R]{\includegraphics[height=1.8cm]{logoUco}} %Logo Universidad
      \renewcommand{\headrulewidth}{0.4pt}
      \renewcommand{\footrulewidth}{0.4pt}
      \setlength{\headheight}{65pt}
    }
    \fancypagestyle{indice}{
      \fancyfoot{}
      \fancyfoot[L]{\AutorTFG} %Autor
      \fancyfoot[R]{- \thepage \ -} %Asignatura
      \fancyhead{}
      \fancyhead[L]{Trabajo Fin de Grado} %Nombre del documento
      \fancyhead[R]{\includegraphics[height=1.8cm]{logoUco}} %Logo
      \renewcommand{\headrulewidth}{0.4pt}
      \renewcommand{\footrulewidth}{0.4pt}
      \setlength{\headheight}{65pt}
      }
      
    \fancypagestyle{plain}{% % <-- this is new
     \fancyfoot{}
      \fancyfoot[L]{Trabajo Fin de Grado} %Autor
      \fancyfoot[C]{} %ecc
      \fancyfoot[R]{- \thepage \ -}
      \fancyhead{}
      \fancyhead[L]{\leftmark} %Nombre del documento
      \fancyhead[R]{\includegraphics[height=1.8cm]{logoUco}} %Logo Universidad
      \renewcommand{\headrulewidth}{0.4pt}
      \renewcommand{\footrulewidth}{0.4pt}
      \setlength{\headheight}{65pt}
    }
  
\setcounter{secnumdepth}{4}

\pagenumbering{empty}


%%%%%%%%%
% ANEXO %
%%%%%%%%%

% Para que ponga ANEXO en vez de APENDICE
 \addto\captionsspanish{\renewcommand{\appendixname}{Anexo}} %para que ponga anexo y no apéndice.

\makeglossaries
\raggedbottom  % Quita huecos grandes en blanco
\begin{document}

%%%%%%%%%%%
% PORTADA %
%%%%%%%%%%%
\input{Portada/cover}
\newpage
\thispagestyle{empty}   % Página en blanco entre la portada y el comienzo del documento
\mbox{} 
$\ $
\thispagestyle{empty} % para que no se numere esta pagina
%\newpage
%$\ $
%\thispagestyle{empty} % para que no se numere esta pagina

%%%%%%%%%%%%%%%%%%%%%%%%%%
% ANEXO V - Originalidad %
%%%%%%%%%%%%%%%%%%%%%%%%%%
\input{AnexoV/anexo_v}
%\includepdf[pages=-]{AnexoV/AnexoV_DECLARACION_ORIGINALIDAD_TFG.PDF}
\newpage
\thispagestyle{empty}
\mbox{}
\newpage

%%%%%%%%%%%%%%%
% DEDICATORIA %
%%%%%%%%%%%%%%%
\pagenumbering{Roman} % para comenzar la numeracion de paginas en numeros romanos
\pagestyle{indice}
\thispagestyle{indice}
\begin{flushright}
\textit{Dedicado a \\
mi familia}
\end{flushright}

%%%%%%%%%%%%%%%%%%%
% AGRADECIMIENTOS %
%%%%%%%%%%%%%%%%%%%
\input{0_agradecimientos}

%%%%%%%%%%%%%%%%%%%%%%
% RESUMEN Y ABSTRACT %
%%%%%%%%%%%%%%%%%%%%%%
\input{0_resumen_y_abstract}

%%%%%%%%%%%%%%%%%%
% INDICE GENERAL %
%%%%%%%%%%%%%%%%%%
\newpage
\thispagestyle{fancy} 

\tableofcontents

\cleardoublepage
\listoffigures

\cleardoublepage
\listoftables

% Descomentar las siguientes líneas si tienes pensado añadir diagramas de flujo.
\cleardoublepage
\phantomsection
\addcontentsline{toc}{chapter}{Índice de diagramas de flujo}
\listofdiagrams

\cleardoublepage
\phantomsection
\addcontentsline{toc}{chapter}{Lista de Acrónimos}

\include{lista_acronimos}
\printglossary[type=\acronymtype, title={Acrónimos}]

%---------------------%
%	CUERPO DE TEXTO   %
%---------------------%

\afterpage{\null\newpage}

\newpage
\newpage

\pagestyle{cuerpo}
\thispagestyle{cuerpo}

\pagenumbering{arabic}

%%%%%%%%%%%%%%%%%%%%%%%%%%%%%%%%%%%%%%%%%%%%%%%%%%%%%%%%%%%%%%%%%%%%%
%                         ESTRUCTURA CONTENIDO                      %
%%%%%%%%%%%%%%%%%%%%%%%%%%%%%%%%%%%%%%%%%%%%%%%%%%%%%%%%%%%%%%%%%%%%%

%------------------%
%	CAPÍTULOS      %
%------------------%

\input{1_introduccion}
\input{2_objetivos}

% --> poner otros capítulos. Por ejemplo:
    %\input{3_antecedentes}    
    %\input{4_materiales}
    %\input{5_diseno}
    %\input{6_pruebas}

% Comentar en la versión final
\input{EJEMPLOS_LATEX}

\input{7_conclusiones}

%--------------------%
% DOCUMENTOS TÉNICOS %
%--------------------%

\input{presupuesto}



%------------------%
%	BIBLIOGRAFIA   %
%------------------%

%Estilo de la bibliografia
\bibliographystyle{unsrtnat} 
%\bibliographystyle{elsarticle-num}
\bibliography{referencias.bib} 

%------------------%
%      ANEXOS      %
%------------------%

\input{planos}

\appendix

\input{anexo_hojas_de_caracteristicas}
\input{anexo_codigo}
\input{anexo_manual_usuario}

\end{document}
